% Level 2, Week 10 answers -- Project Studio 1
\answersweek{10}{Project Studio 1}{Answers and scoring notes for workbook chapter 10 and homework 10.}

\begin{teachernote}[Scoring project weeks]
Project weeks have no single right answer. Score against the milestone: \textbf{met} (it works),
\textbf{partly} (started, not working yet) or \textbf{not yet}. Praise small working steps. The full
project rubric is in the Project Teacher Guide.
\end{teachernote}

\section{Workbook exercises}
\begin{answer}[Exercise 10.1: Paper prototype and plan check (M1)]
A good plan names the loop, at least one decision, the data and 3 or more functions. Example for an
arcade game:
\begin{outputcode}
set up screen, basket, keys (Left/Right), score = 0, lives = 3, best = load_high_score()
every tick (50 ms):
    move every star in the stars list down
    for each star: if touching(basket, star): add 1 to score, move star to top
    if a star reaches the bottom: lose a life
    if lives == 0: if score > best: save_high_score(score); show GAME OVER
\end{outputcode}
Example for an Art Studio:
\begin{outputcode}
shapes = empty list
repeat until the user types q:
    ask: d = draw, s = save, l = load, q = quit
    d: make a random shape dictionary, add it to shapes, draw it
    s: save_design(shapes, "design.txt")
    l: shapes = load_design("design.txt"), clear screen, draw them all
\end{outputcode}
\textbf{Look for:} the three checklist items checked honestly. A plan with no data yet is fine at M1 if
the team can say where the list, dictionary and file will go.
\end{answer}

\begin{answer}[Exercise 10.2: A fair task board]
Typical functions for a catching game: \texttt{next\_x(x, direction)} → new x; \texttt{touching(x1, y1,
x2, y2, how\_close)} → True/False; \texttt{add\_score(score, points)} → new score;
\texttt{load\_high\_score(filename)} → number; \texttt{save\_high\_score(filename, score)} → nothing (writes a
file); \texttt{draw\_star(t, x, y)} → nothing (draws). \textbf{Fair split:} each student owns at
least one logic function (these get tested next week) and at least one drawing or key function.
Watch for one student holding all the interesting notes.
\end{answer}

\begin{answer}[Exercise 10.3: Build the core loop (M2)]
The arcade skeleton is printed in the lesson plan (week 10, \texttt{mygame.py}). Art Studio teams
build a \texttt{while} menu loop. This version was run with the test turtle and in a real window:
\begin{pycode}[studio.py -- Art Studio core loop (tested)]
import random
import turtle

PALETTE = ["#E69F00", "#56B4E9", "#009E73", "#0072B2", "#D55E00", "#CC79A7"]

def random_shape():
    """Make one shape recipe as a dictionary."""
    return {"x": random.randint(-250, 250), "y": random.randint(-200, 200),
            "size": random.randint(10, 60), "color": random.choice(PALETTE)}

def draw_shape(pen, shape):
    pen.penup()
    pen.goto(shape["x"], shape["y"])
    pen.dot(shape["size"], shape["color"])

def studio():
    pen = turtle.Turtle()
    pen.hideturtle()
    shapes = []                       # the list remembers the design
    while True:                       # the core loop: one turn per choice
        choice = input("d = draw, q = quit: ")
        if choice == "q":
            break
        if choice == "d":
            shape = random_shape()
            shapes.append(shape)
            draw_shape(pen, shape)
    print(f"Your design has {len(shapes)} shapes.")

if __name__ == "__main__":
    studio()
    turtle.done()
\end{pycode}
\textbf{M2 met} when: it runs without an error, something repeats by itself (timer or menu; for a team
using \texttt{art\_studio\_template.py}, a key that draws a new design each time), one key
or menu choice works, start-up code is under the main guard. \textbf{Partly}: it runs but nothing
repeats (usually \texttt{ontimer} is missing at the end of \texttt{game\_loop}) or keys do nothing
(missing \texttt{screen.listen()}, or \texttt{onkey(go\_left(), "Left")} with parentheses).
\end{answer}

\begin{answer}[Exercise 10.4: Partner code review]
Good stars are specific: \enquote{\texttt{next\_x} has a clear name and a docstring},
\enquote{all the speeds are in the SETTINGS dictionary}. Good wishes are kind and doable:
\enquote{the function \texttt{do\_stuff} could have a name that says what it does}. Gently redirect
comments about the person (\enquote{you are slow}) to comments about the code.
\end{answer}

\section{Homework}
\begin{answer}[Core: dev log]
Any honest three-line log is a full answer, including \enquote{I got stuck on line 14, the error
said NameError: name 'player' is not defined}. That is exactly the kind of note that helps next week.
\end{answer}

\begin{answer}[Challenge: tests on paper]
For \texttt{next\_x} with a step of 20 and edges at $\pm 280$: \texttt{next\_x(0, 1)} → \texttt{20}
(normal); \texttt{next\_x(280, 1)} → \texttt{280} (edge: stays on screen); \texttt{next\_x(-280, -1)} →
\texttt{-280} (edge); \texttt{next\_x(270, 1)} → \texttt{280} (edge: stops exactly at the wall). Accept any
function whose examples match what the code should do.
\end{answer}

\begin{answer}[Creative: title screen sketch]
Look for a name, the keys explained in words, palette colors and at least one non-color difference
(shape, size, label). Nickname only.
\end{answer}
